\section{Independent-training prewhitening}
\label{app:whitening}

Let $d=Cn$ be fixed and suppose
\[
\operatorname{vec}(\bar Y)=\operatorname{vec}(M)+m^{-1/2}\xi,
\qquad \xi\sim N_d(0,\Sigma),
\]
where the eigenvalues of $\Sigma$ lie in a fixed interval
$[\underline\sigma^2,\overline\sigma^2]$. Independently observe $K$ i.i.d.
$N_d(0,\Sigma)$ noise-training vectors and let $\widehat\Sigma_K$ be their
sample covariance. All fits and projections below use the estimated
inner product
$\langle x,y\rangle_{\widehat\Sigma^{-1}}=x^T\widehat\Sigma^{-1}y$.

\begin{theorem}[plug-in covariance phase score]
\label{thm:whitening}
Assume $K\to\infty$, the null parameters remain in the same compact regular
set as Theorem~\ref{thm:score}, and the eigenvalue bounds above hold uniformly.
Construct the fitted null tangent $\widehat{\mathcal T}$ and project the raw
collision directions
\[
\{u(f'_{\widehat\theta})^T:u\in\mathbb R^C\}
+\operatorname{span}\{\widehat s(f''_{\widehat\theta})^T\}
\]
off $\widehat{\mathcal T}$ in the
$\widehat\Sigma^{-1}$ metric.  Let $\widehat{\mathcal E}_\Sigma$ be the
resulting $C$-dimensional space.  Then the squared norm of the whitened fitted
residual projected onto $\widehat{\mathcal E}_\Sigma$, multiplied by $m$; call
this statistic $T_{m,K}$. Under compact collision-local alternatives covered
by Theorem~\ref{thm:distance}, its distribution is uniformly $o(1)$ in
Kolmogorov distance from $\chi_C^2(\rho_{m,\Sigma}^2)$, where
\[
\rho_{m,\Sigma}^2=
m\|M-\Pi_{\mathcal M_1}^{\Sigma^{-1}}M\|_{\Sigma^{-1}}^2+o(1).
\]
Under the null $\rho_{m,\Sigma}^2=0$, so $T_{m,K}\Rightarrow\chi_C^2$; if
$\rho_{m,\Sigma}^2\to\rho_\Sigma^2$, then
$T_{m,K}\Rightarrow\chi_C^2(\rho_\Sigma^2)$.
No restriction on the relative rates of $K$ and $m$ is needed beyond
$K\to\infty$ in fixed dimension.
\end{theorem}

\begin{proof}
Fixed-dimensional Gaussian covariance estimation and the uniform eigenvalue
bounds imply
$\|\widehat\Sigma_K-\Sigma\|_{\rm op}=O_p(K^{-1/2})$ and, by functional
calculus,
$\|\widehat\Sigma_K^{-1/2}-\Sigma^{-1/2}\|_{\rm op}=O_p(K^{-1/2})$.
Independence makes these perturbations external to the signal central
sequence.  Uniform nonsingularity of the whitened null tangent Gram matrix
then gives the usual root-$m$ expansion of the generalized least-squares null
fit.  QR/SVD perturbation of the tangent and collision spaces is
$O_p(m^{-1/2}+K^{-1/2})$.  Acting on the root-$m$ fitted residual, which is
$O_p(1)$ on compact local sets, this perturbation is $o_p(1)$.  In an
orthonormal frame of the population whitened collision space, the resulting
coordinates therefore equal $h_\Sigma+Z+o_p(1)$ with
$Z\sim N_C(0,I_C)$.  Squaring proves both limits.  Since the two perturbations
enter additively after root-$m$ scaling, $K/m$ need not converge to a positive
constant.
\end{proof}

For ordered rates, nonnegative component amplitudes, and a strictly positive
merged-amplitude vector, the physical local parameter set is a half-space
before whitening. Its image under the
nonsingular coordinate map to $\widehat{\mathcal E}_\Sigma$ is again a
half-space.  Projection onto its consistently estimated inward normal yields
the same equal-weight
$\frac12\chi^2_{C-1}+\frac12\chi^2_C$ null limit.  The result covers general
cross-channel and across-time covariance, provided the training vectors are
independent and generated under the same covariance law as the signal noise.
It does not justify covariance estimated from the tested fitted residual.

The theorem is fixed-dimensional and asymptotic in the independent training
size.  If the full covariance is estimated with $K$ only slightly larger than
$d$, neither the chi-square critical value nor a scalar inverse-Wishart
centering correction is exact because the estimated metric also changes the
fitted tangent and collision spaces.  Structured covariance estimation or a
finite-training calibration is therefore required in that regime.
